\writeSectionFrame{\presentationTitle}{Das Problem}
\begin{frame}{Ohne Adapter}
	\begin{center}
 		\includegraphics[scale=0.3]{\BILDERPFAD/problem1.png}
 	\end{center}
\end{frame}

\note{
	Wir möchten Daten von einem Wetterdienst abrufen und für die Benutzer auf
	einer wunderschönen Webseite aufbereiten, anzeigen und ergänzen. Jetzt erhalten wir eine Anfrage
	von einem Forschungsinstitut, ob wir die Daten nicht an sie weitergeben könnten. Das
	wäre an sich kein Problem, aber leider arbeiten unsere Klassen und der Wetterdienst mit
	xml, während das Forschungsinstitut json-Dateien erwartet. Andere Forschungsinstitute melden sich
	dann auch. Einige können allerdings auch mit xml-Dateien umgehen. Solche Fälle, in denen
	bestimmte Dienste nicht kompatibel zueinander sind, findet man in der IT sehr häufig. 
}

\begin{frame}{Mit Adapter}
	\begin{center}
 		\includegraphics[scale=0.3]{\BILDERPFAD/problem2.png}
 	\end{center}
\end{frame}

\note{
 	Hier bietet ein Adapter eine Lösung. Im Idealfall fügt er sich nahtlos in die Struktur ein,
 	so dass die beteiligten Klassen gar nicht bemerken, dass ein Adapter dazwischen geschaltet ist. 
}